\documentclass[12pt]{article}
\usepackage{amsmath,amssymb,amsfonts}
\usepackage[margin=1in]{geometry}

\begin{document}

\textbf{Chapter 6, Problem 32.} We have seen how the function $e\cot ez$ may be used to sum certain series. But it is useless if the series is an alternating one. Show how alternating series may be summed by using the function $e\csc(ez)$ in place of $e\cot(ez)$. Use the general result to prove that
\[
\sum_{n=1}^{\infty}(-1)^{n+1}\frac{1}{n^2} = \frac{\pi^2}{12}.
\]

\bigskip
\textbf{Solution.}

The function $\pi\csc(\pi z)$ has simple poles at $z = n$ (all integers) with residues $(-1)^n$:
\[
\text{Res}_{z=n}\pi\csc(\pi z) = \frac{\pi}{(\pi\sin(\pi z))'|_{z=n}} \cdot \pi = \frac{1}{\cos(n\pi)} = (-1)^n.
\]

Wait, more carefully: $\csc(\pi z) = 1/\sin(\pi z)$ has poles at $z = n$ with $\text{Res}_{z=n} = \frac{1}{\pi\cos(n\pi)} = \frac{(-1)^n}{\pi}$. So $\pi\csc(\pi z)$ has residue $(-1)^n$ at $z = n$.

Consider the contour integral of $\pi\csc(\pi z)\cdot g(z)$ around a large square contour $C_N$ enclosing $z = -N, \ldots, N$:
\[
\frac{1}{2\pi i}\oint_{C_N}\pi\csc(\pi z)g(z)\,dz = \sum_{n=-N}^{N}(-1)^n g(n) + \sum_{\text{poles of }g}\text{Res}[\pi\csc(\pi z)g(z)].
\]

For $g(z) = 1/z^2$: poles of $g$ at $z = 0$ (order 2). The combined function $\pi\csc(\pi z)/z^2$ near $z = 0$:

$\pi\csc(\pi z) = \frac{1}{z} + \frac{\pi^2 z}{6} + \cdots$, so $\frac{\pi\csc(\pi z)}{z^2} = \frac{1}{z^3} + \frac{\pi^2}{6z} + \cdots$

The residue of the combined function at $z = 0$ is $\pi^2/6$.

As $N \to \infty$, the contour integral vanishes (since $|\csc(\pi z)|$ is bounded on the square and $1/z^2 \to 0$). The sum of all residues is zero:
\[
0 = \sum_{n\neq 0}(-1)^n \frac{1}{n^2} + \frac{\pi^2}{6}.
\]

\[
\sum_{n\neq 0}(-1)^n\frac{1}{n^2} = -\frac{\pi^2}{6}.
\]

\[
\sum_{n=1}^{\infty}\frac{(-1)^n}{n^2} + \sum_{n=1}^{\infty}\frac{(-1)^{-n}}{n^2} = 2\sum_{n=1}^{\infty}\frac{(-1)^n}{n^2} = -\frac{\pi^2}{6}.
\]

\[
\sum_{n=1}^{\infty}\frac{(-1)^n}{n^2} = -\frac{\pi^2}{12}.
\]

Therefore:
\[
\boxed{\sum_{n=1}^{\infty}\frac{(-1)^{n+1}}{n^2} = \frac{\pi^2}{12}.}
\]

\end{document}
